\section{Method}
\label{sec:method}

\subsection{Problem Setup}
\label{sec:setup}

Consider a constellation of $N$ satellites, each parameterized by a two-line element (TLE) set.
For constellation geometry optimization at a fixed altitude, the free parameters per satellite are the inclination $\iota$, right ascension of the ascending node $\Omega$, and mean anomaly $\mathcal{M}$.
We fix the remaining TLE parameters (mean motion $n$, eccentricity $e \approx 0$, drag term $B^*$, and secular rates $\dot{n}$, $\ddot{n}$).

Given a grid of $M$ ground target points $\mathcal{G} = \{\vect{g}_1, \ldots, \vect{g}_M\}$ on the Earth's surface and a time horizon $[0, T]$ discretized into $K$ steps with spacing $\delta t = T/(K-1)$, we seek to maximize coverage and minimize the worst-case revisit gap.

\paragraph{Propagation.}
Each satellite's TLE elements are propagated through dSGP4~\citep{acciarini2024dsgp4} to produce position vectors $\vect{r}_i(t_k)$ in the True Equator Mean Equinox (TEME) reference frame.
These are rotated to the Earth-Centered Earth-Fixed (ECEF) frame via the Greenwich Mean Sidereal Time angle $\theta_{\text{GMST}}(t_k)$:
\begin{equation}
    \vect{r}_i^{\text{ECEF}}(t_k) = \mat{R}_z\bigl(\theta_{\text{GMST}}(t_k)\bigr) \, \vect{r}_i^{\text{TEME}}(t_k),
    \label{eq:teme_to_ecef}
\end{equation}
where $\mat{R}_z(\theta)$ is the standard rotation matrix about the $z$-axis.
This single-axis rotation is an approximation: the full TEME-to-ECEF transformation requires accounting for precession, nutation, and polar motion~\citep{vallado2006revisiting}.
For the \SI{24}{\hour} propagation horizons considered here, the precession-induced error is on the order of arcseconds, corresponding to sub-kilometre position offsets---negligible relative to the $\sim$\SI{2200}{\kilo\metre} ground footprint of a satellite at \SI{550}{\kilo\metre} altitude with \SI{10}{\degree} minimum elevation.
For longer horizons or higher-precision applications, the full IAU-2000/2006 transformation should be used.

\paragraph{Elevation angle.}
The elevation angle from ground point $\vect{g}_j$ to satellite $i$ at time $t_k$ is
\begin{equation}
    \alpha_{ij}(t_k) = \arcsin\left(\frac{(\vect{r}_i^{\text{ECEF}}(t_k) - \vect{g}_j) \cdot \hat{\vect{g}}_j}{\|\vect{r}_i^{\text{ECEF}}(t_k) - \vect{g}_j\|}\right),
    \label{eq:elevation}
\end{equation}
where $\hat{\vect{g}}_j = \vect{g}_j / \|\vect{g}_j\|$ is the local vertical unit vector at the ground point.
A satellite is conventionally visible when $\alpha_{ij} \geq \alpha_{\min}$, with $\alpha_{\min} = \SI{10}{\degree}$.

The challenge is that the visibility indicator $\mathbf{1}[\alpha_{ij}(t_k) \geq \alpha_{\min}]$, the Boolean OR over satellites, the revisit gap computation, and the worst-case maximum are all non-differentiable.
We address each in turn.

\subsection{Relaxation 1: Soft Visibility}
\label{sec:soft_visibility}

We replace the hard step function with a sigmoid (\figref{fig:relaxations}a):
\begin{equation}
    c_{ij}(t_k) = \sigma\!\left(\frac{\alpha_{ij}(t_k) - \alpha_{\min}}{\tau}\right), \quad \sigma(x) = \frac{1}{1 + e^{-x}},
    \label{eq:soft_visibility}
\end{equation}
where $\tau > 0$ controls the sharpness.
As $\tau \to 0$ this recovers the hard indicator; we use $\tau = \SI{2}{\degree}$, which provides a smooth transition zone of approximately $\pm\SI{6}{\degree}$ around the threshold while keeping the approximation tight.
This is analogous to the temperature-scaled softmax used in Gumbel-Softmax relaxations of discrete variables~\citep{jang2017categorical}.

\subsection{Relaxation 2: Noisy-OR Coverage Aggregation}
\label{sec:noisy_or}

For a ground point $\vect{g}_j$ at time $t_k$, we need to determine whether \emph{any} satellite provides coverage---a Boolean OR.
Treating each $c_{ij}(t_k)$ as an independent detection probability, the noisy-OR~\citep{pearl1988probabilistic} gives (\figref{fig:relaxations}b):
\begin{equation}
    C_j(t_k) = 1 - \prod_{i=1}^{N} \bigl(1 - c_{ij}(t_k)\bigr).
    \label{eq:noisy_or}
\end{equation}
This is smooth, monotonically increasing in each $c_{ij}$, and recovers the hard OR when all inputs are in $\{0, 1\}$.
It also naturally handles overlapping coverage: adding a satellite that already covers a point produces diminishing marginal gain, preventing the optimizer from clustering satellites.

The \emph{coverage fraction} over the full space-time grid is
\begin{equation}
    \bar{C} = \frac{1}{KM} \sum_{k=1}^{K} \sum_{j=1}^{M} C_j(t_k).
    \label{eq:coverage_fraction}
\end{equation}

\subsection{Relaxation 3: Leaky Integrator for Revisit Gaps}
\label{sec:leaky_integrator}

The revisit gap at a ground point is the time elapsed since the last coverage event.
Computing this requires identifying discrete transitions, which is non-differentiable.
We approximate it with a leaky integrator recurrence~\citep{gerstner2002spiking} (\figref{fig:relaxations}c):
\begin{equation}
    \Delta_j(t_k) = \bigl(\Delta_j(t_{k-1}) + \delta t\bigr) \cdot \bigl(1 - C_j(t_k)\bigr), \quad \Delta_j(t_0) = 0.
    \label{eq:leaky_integrator}
\end{equation}
When the point is covered ($C_j \approx 1$), the accumulated gap is multiplied by $\approx 0$, resetting it.
When the point is not covered ($C_j \approx 0$), the gap grows by $\delta t$.
This recurrence is fully differentiable and its maximum over time approximates the true worst-case revisit gap at each ground point.

\subsection{Relaxation 4: LogSumExp Soft-Maximum}
\label{sec:logsumexp}

The worst-case revisit gap at ground point $\vect{g}_j$ is $\max_k \Delta_j(t_k)$, which has zero gradients almost everywhere.
We replace it with the LogSumExp smooth maximum~\citep{boyd2004convex} (\figref{fig:relaxations}d):
\begin{equation}
    \widetilde{\Delta}_j^{\max} = \tau_r \cdot \log\!\left(\sum_{k=1}^{K} \exp\!\left(\frac{\Delta_j(t_k)}{\tau_r}\right)\right),
    \label{eq:logsumexp}
\end{equation}
where $\tau_r = \SI{30}{\minute}$.
This satisfies $\max_k \Delta_j(t_k) \leq \widetilde{\Delta}_j^{\max} \leq \max_k \Delta_j(t_k) + \tau_r \log K$, providing a controllable upper bound.
The mean worst-case revisit across all ground points is
\begin{equation}
    \bar{\Delta}^{\max} = \frac{1}{M} \sum_{j=1}^{M} \widetilde{\Delta}_j^{\max}.
    \label{eq:mean_revisit}
\end{equation}

\begin{figure}[t]
    \centering
    \begin{tikzpicture}
\begin{groupplot}[
    group style={
        group size=4 by 1,
        horizontal sep=1.4cm,
    },
    width=4.2cm, height=3.5cm,
    grid=major,
    grid style={gray!20},
    tick label style={font=\scriptsize},
    label style={font=\small},
    title style={font=\small\bfseries, yshift=-2pt},
]

% (a) Soft sigmoid
\nextgroupplot[
    title={(a) Soft Visibility},
    xlabel={$\alpha - \alpha_{\min}$ [\si{\degree}]},
    ylabel={$c_{ij}$},
    xmin=-15, xmax=15,
    ymin=-0.05, ymax=1.1,
    domain=-15:15,
    samples=100,
]
% Hard step
\addplot[blue, very thick, dashed] coordinates {(-15,0) (0,0) (0,1) (15,1)};
% Soft sigmoid tau=2
\addplot[red, very thick] {1/(1+exp(-x/2))};
\node[font=\scriptsize] at (axis cs:7,0.3) {$\tau=2$};

% (b) Noisy-OR
\nextgroupplot[
    title={(b) Noisy-OR},
    xlabel={$c_1$},
    ylabel={$C = 1 - (1-c_1)(1-c_2)$},
    xmin=0, xmax=1,
    ymin=-0.05, ymax=1.1,
    domain=0:1,
    samples=50,
]
\addplot[red, very thick] {1-(1-x)*(1-0.0)};
\addplot[orange, very thick] {1-(1-x)*(1-0.3)};
\addplot[purple, very thick] {1-(1-x)*(1-0.7)};
\node[font=\scriptsize, anchor=west] at (axis cs:0.05,0.15) {$c_2\!=\!0$};
\node[font=\scriptsize, anchor=west] at (axis cs:0.05,0.42) {$c_2\!=\!0.3$};
\node[font=\scriptsize, anchor=west] at (axis cs:0.05,0.78) {$c_2\!=\!0.7$};

% (c) Leaky integrator
\nextgroupplot[
    title={(c) Leaky Integrator},
    xlabel={Time step},
    ylabel={Gap $\Delta_j$ (min)},
    xmin=0, xmax=20,
    ymin=-5, ymax=85,
]
% Simulated gap signal: grows linearly, resets at coverage events
\addplot[red, very thick] coordinates {
    (0,0) (1,7) (2,14) (3,21) (4,28) (5,35)
    (6,0) (7,7) (8,14) (9,21) (10,28) (11,35) (12,42)
    (13,0) (14,7) (15,14) (16,21) (17,28) (18,35) (19,42) (20,49)
};
% Coverage events
\addplot[blue, only marks, mark=triangle*, mark size=2.5pt] coordinates {(6,0) (13,0)};
\node[font=\scriptsize, blue] at (axis cs:9.5,65) {coverage};

% (d) LogSumExp
\nextgroupplot[
    title={(d) LogSumExp Soft-Max},
    xlabel={Temperature $\tau_r$},
    ylabel={$\widetilde{\Delta}^{\max}$},
    xmin=0.5, xmax=50,
    ymin=40, ymax=110,
    domain=1:50,
    samples=50,
]
% For values [20, 45, 30], true max = 45
\addplot[red, very thick] {x * ln(exp(20/x) + exp(45/x) + exp(30/x))};
\addplot[blue, very thick, dashed] coordinates {(0.5,45) (50,45)};
\node[font=\scriptsize, blue] at (axis cs:35,42) {true max};
\node[font=\scriptsize, red, anchor=west] at (axis cs:25,85) {LSE};

\end{groupplot}
\end{tikzpicture}

    \caption{The four continuous relaxations.  \textbf{(a)} Soft sigmoid replaces the hard visibility threshold.  \textbf{(b)} Noisy-OR aggregates coverage across satellites.  \textbf{(c)} Leaky integrator tracks the revisit gap as a differentiable recurrence.  \textbf{(d)} LogSumExp approximates the worst-case (maximum) revisit gap.}
    \label{fig:relaxations}
\end{figure}

\subsection{Combined Objective}
\label{sec:objective}

The total loss combines coverage maximization with revisit minimization:
\begin{equation}
    \mathcal{L}(\vect{x}) = -\bar{C} + \lambda \, \bar{\Delta}^{\max},
    \label{eq:loss}
\end{equation}
where $\vect{x} = \{\iota_i, \Omega_i, \mathcal{M}_i\}_{i=1}^{N}$ collects all free orbital parameters and $\lambda$ is a weighting coefficient.
For non-uniform objectives, the sums over ground points can be replaced with weighted sums over a discrete target set (\secref{sec:weighted}).

\subsection{Tightness of the Composed Relaxation}
\label{sec:tightness}

Each of the four relaxations is individually tight: the soft sigmoid recovers the hard step as $\tau \to 0$; the noisy-OR is exact when its inputs are binary; the leaky integrator is exact when coverage is in $\{0, 1\}$; and the LogSumExp satisfies $\max_k x_k \leq \text{LSE} \leq \max_k x_k + \tau_r \log K$.
However, the \emph{composition} of these relaxations does not inherit these guarantees, because each relaxation receives the soft output of the previous one rather than the binary values for which it was designed.

The interaction introduces two competing biases:

\paragraph{Phantom coverage (gap shortening).}
The soft sigmoid assigns nonzero coverage $c_{ij} > 0$ to satellites slightly \emph{below} the visibility threshold ($\alpha_{ij} < \alpha_{\min}$).
This phantom coverage propagates through the noisy-OR and partially resets the leaky integrator, producing gaps that are \emph{shorter} than the true discrete gaps.
For example, a satellite at $\alpha = \alpha_{\min} - 2\tau$ contributes $c_{ij} \approx 0.12$, which the leaky integrator treats as a partial coverage event.

\paragraph{Incomplete reset (gap lengthening).}
Conversely, a satellite well above the threshold produces $c_{ij} \approx 0.99$, not exactly 1.
The leaky integrator retains a fraction $(1 - C_j) \approx 0.01$ of the accumulated gap at each coverage event.
Over time, this residual prevents the gap from fully resetting, producing relaxed gaps that are \emph{longer} than the true gaps during sustained coverage.
The LogSumExp adds a further upward bias of at most $\tau_r \log K$.

These effects compete: phantom coverage shortens gaps near the visibility boundary, while incomplete resets and the LogSumExp overestimate lengthen them.
The net bias depends on the specific geometry and is neither a guaranteed upper nor lower bound on the true objective.

For optimization, the critical question is not whether the relaxed loss equals the true loss, but whether the two have similar \emph{optima}.
We validate this empirically in \secref{sec:relaxation_quality} by computing the hard (discrete) coverage and revisit metrics at the optimizer's solution and comparing them to the relaxed values.
The relaxation temperatures $\tau$ and $\tau_r$ control the trade-off: smaller values tighten the approximation but sharpen the gradients, potentially destabilizing optimization.
A temperature annealing schedule---starting with large temperatures for a smooth landscape and progressively tightening---could mitigate this trade-off, though we leave this to future work.

\subsection{Optimization}
\label{sec:optimization}

The full computational graph is:
\begin{equation}
    \underbrace{\vect{x}_i}_{\text{TLE}} \xrightarrow{\text{dSGP4}} \underbrace{\vect{r}_i(t)}_{\text{TEME}} \xrightarrow{\text{GMST}} \underbrace{\vect{r}_i^{\text{ECEF}}(t)}_{\text{ECEF}} \xrightarrow{\text{Eq.~\ref{eq:elevation}}} \alpha_{ij}(t) \xrightarrow{\text{Eqs.~\ref{eq:soft_visibility}--\ref{eq:mean_revisit}}} \mathcal{L}.
    \label{eq:pipeline}
\end{equation}
Every operation is implemented in PyTorch~\citep{paszke2019pytorch}, so $\nabla_{\vect{x}_i} \mathcal{L}$ is computed via reverse-mode automatic differentiation.

A practical complication is that $\partial$SGP4's \texttt{initialize\_tle} function rebuilds the computation graph from scratch at each call, producing a new ephemeral tensor for each satellite's elements.
Standard optimizers like \texttt{torch.optim.Adam} track specific tensor objects and would lose their momentum state.
We solve this by maintaining persistent parameter tensors that Adam tracks, and copying gradients from dSGP4's ephemeral tensors after each backward pass:
\begin{equation}
    \vect{p}_i.\text{grad} \leftarrow \text{clip}\bigl(\nabla_{\vect{x}_i^{\text{eph}}} \mathcal{L},\ [-5, 5]\bigr) \odot \vect{m},
    \label{eq:grad_copy}
\end{equation}
where $\vect{p}_i$ is the persistent parameter tensor, $\vect{x}_i^{\text{eph}}$ is dSGP4's ephemeral tensor, and $\vect{m} \in \{0,1\}^9$ is a binary mask that zeroes gradients for fixed parameters (altitude, eccentricity, drag).
Adam then steps on $\vect{p}_i$, and the updated values are written back to the TLE object before the next iteration.

The per-iteration cost is $\mathcal{O}(NKM)$: $N$ satellite propagations over $K$ time steps, each evaluated against $M$ ground points.
